\documentclass[a4paper,12pt,oneside]{scrreprt}
\usepackage[latin1]{inputenc}
\usepackage[T1]{fontenc}
\usepackage{ae,aecompl}
\usepackage[english]{babel}
\usepackage{amsmath}
\usepackage{amssymb}
\usepackage{amsfonts}
\usepackage{amsthm}
\usepackage{graphicx}
\usepackage{wrapfig}
\usepackage{ulem}
\usepackage{cancel}
\usepackage{float}
\usepackage{color}
%\usepackage{titlesec}
\usepackage{geometry}
\geometry{verbose,a4paper,tmargin=25mm,bmargin=25mm,lmargin=15mm,rmargin=25mm}
%\titlelabel{\thetitle.\quad}
\usepackage{tabularx}
\usepackage{booktabs}
\usepackage{paralist}
\usepackage{textcomp}
\usepackage[official]{eurosym}

\renewcommand{\rmdefault}{phv}
\renewcommand{\sfdefault}{phv}

% our OWN imports for our use
\usepackage{listings}
\lstdefinestyle{ourJavaStyle}{
	language=Java,
	%numbers=left,
	%numbersep=8pt,
	stepnumber=1,
	tabsize=2,
	showspaces=false,
	showstringspaces=false,
	basicstyle=\ttfamily\scriptsize,
	keywordstyle=\color{blue}\ttfamily,
	stringstyle=\color{red}\ttfamily,
	commentstyle=\color{green}\ttfamily,
	breaklines=true
}

\newcommand*{\sourcepath}{../code/src/main/java/de/rwth/swc/group10}
\newcommand*{\testpath}{../code/src/test/java/de/rwth/swc/group10}

\renewcommand{\thesubsection}{\thesection.\alph{subsection}}
%\titleformat{\subsection}
%{\normalfont\fontfamily{phv}\fontsize{14}{17}\bfseries}{\thesubsection}{1em}{}

\begin{document}
	
	\begin{tabular}{ccc}
		\begin{large} \textbf{Prof. Lichter} \end{large} &
		
		\begin{minipage}[H]{3.5cm}
			\centering
			\begin{large} OOSC \end{large} \\
			\begin{large} WS 2019/2020 \end{large}
		\end{minipage} &
		
		\begin{minipage}[H]{4cm}
			\includegraphics[keepaspectratio,width=\textwidth,angle=0]{images/swc.png}
		\end{minipage} \\
		Andreas Steffens, Konrad F\"ogen &  &  \\
		& \begin{huge} \textbf{Submission 2} \end{huge}&  \\
		& oosc@swc.rwth-aachen.de &  \\
		& & \\
		% Hier drunter muessen die Daten noch angepasst werden
		Issued: 04.11.2019 &
		Submission: 18.11.2019 &
		Discussion: 21.11.2019 \\
	\end{tabular}
	\newline \newline \newline
	\centering
	Submitted by Group 10
	
	\begin{tabular}{ll}
		Dominik Bittner, & 369202 \\
		Ulfet Cetin, & 391819\\
		Philipp Hochmann, & 356148 \\
		Anar Orujov, & 391825\\
		Ada Slupczynski, & 384147\\
		(sorted on lastname basis)
	\end{tabular}
	
	\setcounter{chapter}{2} % Aktuelles Assigment
	\section{Assignment:}
	
    % TODO: Image of the UML is missing
    
	\begin{table}[h]
		\centering
		\resizebox{\textwidth}{!}
		{%
			\begin{tabular}{|l|l|}
				\hline
				\multicolumn{1}{|c|}{\textbf{Statement}} & \multicolumn{1}{c|}{\textbf{Output or Description of Failure}} \\ \hline
				cow.writeName();                         & output: "Cow"         \\ \hline
				aurochs.writeName();                     & output: "Aurochs"     \\ \hline
				((Material)aurochs).writeName();         & output: "Material"    \\ \hline
				((Aurochs)cow).writeName();              & output: "Aurochs"     \\ \hline
				((Cow)aurochs).writeName();              & output: 
                "Cow"         \\ \hline
				aurochs.writeLatin();                    & output: "Bos primigenius"  \\ \hline
				auerochse.writeDescription();              & fail:
                variable auerochse is not defined.      \\ \hline
				((Aurochs)material).writeName();         & fail: we cannot limit to subclass, variable material is not a Aurochs.
				                                                                 \\ \hline
			\end{tabular}
		}
	\end{table}
	
	\section{Assignment:}
	\section{Assignment:}
	\section{Assignment:}
	\section{Assignment:}
	\section{Assignment:}
\end{document}